\documentclass[11pt,a4paper]{article}
\usepackage[T1]{fontenc}
\usepackage[utf8]{inputenc}
\usepackage[french]{babel}
\usepackage{amsmath,amssymb,amsthm}
\usepackage{geometry}
\geometry{margin=2.5cm}
\setlength{\parindent}{0pt}
\newcommand{\normal}{\mathrel{\triangleleft}}
\newcommand{\semi}{\rtimes}

\title{TD : Groupes quotients et sous-groupes normaux}
\author{Vincent Boulard}
\date{October 23, 2023}

\begin{document}
\maketitle

\subsection*{Exercice 1 (sous-groupe normal, quotient, rappels)}
Soit $G$ un groupe et $H\subset G$ un sous-groupe de $G$, on dit que $H$ est distingué (ou \emph{normal}) dans $G$ si
\[
  \forall g\in G,\ \forall h\in H,\quad ghg^{-1}\in H.
\]
Et dans ce cas on note $H\normal G$.
\begin{enumerate}
  \item Montrer que la définition précédente équivaut à $\forall g\in G,\ gH=Hg$.
  \item On note $G/H$ l'ensemble des $gH$ avec $g\in G$. Montrer que $G/H$ est l'ensemble des classes d'équivalence pour la relation suivante $\sim_H$ définie par
  \[
    \forall g,g'\in G,\quad g\sim_H g' \Longleftrightarrow g^{-1}g'\in H.
  \]
  Ainsi on notera $\bar g=gH$ la classe d'équivalence de $g$ pour cette relation. Montrer alors que si $H\normal G$ l'application $(\bar g,\bar g')\mapsto \overline{gg'}$ est bien définie\footnote{C'est une application de $G/H\times G/H$ dans $G/H$ or elle est définie via des représentants. Ainsi il faut vérifier que quand $H$ est normal dans $G$, la définition de celle-ci dépend uniquement des classes de $g$ et $g'$.}, et que munie de celle-ci $G/H$ est un groupe. On dit que c'est le \emph{groupe quotient} de $G$ par $H$.
\end{enumerate}

\subsection*{Exercice 2 (questions en vrac)}
\begin{enumerate}
  \item Soit $k$ un corps, montrer que $SL_n(k)\normal GL_n(k)$ et que $GL_n(k)/SL_n(k)\simeq k^*$.\footnote{$G_1\simeq G_2$ signifie que $G_1$ est isomorphe à $G_2$.}
  \item On note $S_n$ le groupe symétrique de degré $n$. Soit $H\normal S_n$ tel que $H$ contienne une transposition, montrer alors que $H=S_n$.
\end{enumerate}

\subsection*{Exercice 3 (centre)}
Soit $G$ un groupe, on appelle centre de $G$ la partie suivante
\[
  Z(G)=\{x\in G\mid \forall g\in G,\ gx=xg\}.
\]
Et pour tout $g\in G$ on pose $f_g:x\mapsto gxg^{-1}$ et on note $\operatorname{Int}(G)=\{f_g,\ g\in G\}$.
\begin{enumerate}
  \item Montrer que $\operatorname{Int}(G)$ est un groupe pour la composition et que $Z(G)$ est un sous-groupe normal de $G$.
  \item Montrer que $G/Z(G)\simeq \operatorname{Int}(G)$.
  \item Montrer que si $G/Z(G)$ est monogène (\emph{i.e. engendré par un élément}) alors $G$ est abélien.\footnote{Cette question est indépendante du lien entre les automorphismes intérieurs et $Z(G)$.}
\end{enumerate}

\subsection*{Exercice 4 (décomposition directe)}
Soit $G$ un groupe et $K\normal G$ et $H\normal G$ tels que $HK=G$ et $H\cap K=\{e\}$. Montrer que $G\simeq H\times K$.

\subsection*{Exercice 5 (groupe dérivé)}
Soit $G$ un groupe. Soit $x,y\in G$, on appelle commutateur de $x$ et $y$ l'élément suivant $[x,y]=xyx^{-1}y^{-1}$ (il \emph{mesure le défaut de commutation car $[x,y]=e\Longleftrightarrow xy=yx$}), puis le sous-groupe $D(G)$ engendré par tous les commutateurs est appelé \emph{groupe dérivé} de $G$.
\begin{enumerate}
  \item Montrer que $D(G)\normal G$.
  \item Soit $H\normal G$, montrer que $G/H$ est abélien si et seulement si $D(G)\subset H$, autrement dit, $D(G)$ est le plus petit sous-groupe normal tel que $G/D(G)$ soit abélien.
\end{enumerate}

\subsection*{Exercice 6 (groupe résoluble, Galois ?)}
Un groupe $G$ est dit résoluble s'il existe $G_0,\ldots,G_n$ des sous-groupes de $G$ tels que
\[
  \{e\}=G_0\normal G_1\normal\cdots\normal G_{n-1}\normal G_n=G,
\]
et tels que $G_{i+1}/G_i$ soit abélien.\footnote{Intuitivement, un groupe résoluble est un groupe ``pas trop loin'' d'un groupe abélien.}
\begin{enumerate}
  \item Montrer que si $G$ est abélien alors il est résoluble.
  \item On pose $D^0(G)=G$ et pour tout $n\in\mathbb N$, $D^{n+1}(G)=D(D^n(G))$. Montrer via l'exercice précédent que $G$ est résoluble si et seulement si la suite $(D^n(G))_{n\in\mathbb N}$ stationne à $\{e\}$.
  \item On se donne un sous-groupe $H$ de $G$.
  \begin{enumerate}
    \item Montrer que si $G$ est résoluble alors $H$ aussi.
    \item Montrer que si $H\normal G$ et que $G$ est résoluble, alors $G/H$ l'est aussi.
    \item Montrer que si $H\normal G$ et que $H$ et $G/H$ sont résolubles, alors $G$ est résoluble.
  \end{enumerate}
\end{enumerate}

\subsection*{Exercice 7 (exemples de groupes résolubles)}
\begin{enumerate}
  \item Montrer que $S_1,S_2$ et $S_3$ sont résolubles.
  \item On note $D_n=\{e,r,r^2,\ldots,r^{n-1},s,rs,\ldots,r^{n-1}s\}$ le groupe diédral d'ordre $2n$\footnote{Le groupe de symétrie d'un $n$-gone régulier centré en $0$ dans le plan, $r$ est la rotation directe d'angle $\frac{2\pi}{n}$ et $s$ la symétrie axiale par rapport à la droite passant par le centre et $1$.}. Montrer que $D_n$ est résoluble (indication : montrer que $D(D_n)$ est engendré par $r^2$).
\end{enumerate}

\subsection*{Exercice 8 (groupe simple)}
On dit qu'un groupe $G$ est simple s'il admet uniquement deux sous-groupes distingués distincts.\footnote{Ainsi on remarque que $\{e\}$ n'est pas simple.}
\begin{enumerate}
  \item Quels sont les groupes abéliens simples ?
  \item Soit $G$ un groupe simple et résoluble, montrer alors qu'il existe un nombre premier $p$ tel que $G\simeq C_p$.\footnote{$C_p$ est le groupe cyclique d'ordre $p$, ainsi il est isomorphe à $\mathbb Z/p\mathbb Z$.}
\end{enumerate}

\subsection*{Exercice 9 (produit semi-direct)}
Soit $N$ et $H$ deux groupes, $f:H\to\operatorname{Aut}(N)$ un morphisme de $H$ dans le groupe des automorphismes de $N$. On munit le produit $N\times H$ de la loi suivante
\[
  ((n,h),(n',h'))\mapsto (n f(h)(n'),hh').
\]
On dit alors que $N\times H$ est le produit semi-direct de $N$ par $H$ suivant $f$ et on note ce groupe $N\semi H$.\footnote{En toute rigueur on devrait souligner davantage le rôle de $f$ pour définir ce groupe, en notant par exemple $\semi_f$ au lieu de $\semi$, mais en pratique on ne le fait pas.}
\begin{enumerate}
  \item Vérifier que la loi définie respecte les axiomes d'une loi de groupe.
  \item On considère les sous-groupes $H=\{e,s\}$ et $N=\{e,r,\ldots,r^{n-1}\}$ de $D_n$, on pose alors $f:H\to\operatorname{Aut}(N)$ tel que $f(e)(r^k)=r^k$ et $f(s)(r^k)=r^{-k}$. En déduire que $D_n\simeq N\semi H$ et donc que $D_n\simeq C_n\semi C_2$.
  \item Soit $V$ un espace vectoriel, on note $GA(V)$ le groupe des automorphismes affines de $V$\footnote{C'est-à-dire les automorphismes de $V$ composés par une translation à gauche.}, montrer alors que $GA(V)=V\semi GL(V)$.
\end{enumerate}

\subsection*{Exercice 10 (décomposition semi-directe)}
Soit $G$ un groupe, $H$ un sous-groupe de $G$ et $N\normal G$ avec $NH=G$ et $N\cap H=\{e\}$. Montrer alors que $G\simeq N\semi H$, comparer avec l'exercice 4.

\emph{Indication : pour définir $f$, penser aux automorphismes intérieurs.}

\subsection*{Exercice 11 (théorème 5/8)}
Soit $G$ un groupe fini de cardinal $n$, on note $z$ le cardinal de son centre. On veut montrer que si l'on tire au hasard deux éléments $g_1$ et $g_2$ de $G$ (\emph{i.e. selon une loi uniforme sur $G^2$}), la probabilité qu'ils commutent est toujours inférieure ou égale à $\frac58$.
\begin{enumerate}
  \item Montrer que le nombre de couples commutant est inférieur à $\frac{zn}{2}+\frac{n^2}{2}$.
  \item En utilisant la dernière question de l'exercice 3, montrer que $z\leq \frac n4$.
  \item Conclure.
\end{enumerate}

\end{document}
